% =====================================================================
% Lab Séance 1 — ADTs et Interfaces Java
% Structures de Données et Algorithmes Avancés — M1 — ESP/UCAD
% Dr. El Hadji Bassirou TOURÉ — 2025–2026
% Compilation : pdflatex (2 passes)
% =====================================================================
\documentclass[11pt,a4paper]{article}

\usepackage[utf8]{inputenc}
\usepackage[T1]{fontenc}
\usepackage{lmodern}
\usepackage[french]{babel}

\usepackage[margin=2.5cm]{geometry}
\usepackage{amsmath,amssymb,amsthm}
\usepackage{graphicx}
\graphicspath{{figures/}}
\usepackage{xcolor}
\usepackage{booktabs}
\usepackage{array}
\usepackage{enumitem}
\usepackage{fancyhdr}
\usepackage{longtable}
\usepackage{multirow}
\usepackage{pifont}
\usepackage{listings}
\usepackage{tcolorbox}
\tcbuselibrary{theorems,breakable,skins}
\usepackage{hyperref}

\hypersetup{colorlinks=true,linkcolor=blue!60!black,urlcolor=blue!60!black,
            citecolor=blue!60!black}

\newcommand{\cmark}{\ding{51}}
\newcommand{\xmark}{\ding{55}}
\newcommand{\bigO}{\mathcal{O}}
\newcommand{\bigOmega}{\Omega}
\newcommand{\bigTheta}{\Theta}
\newcommand{\lig}[1]{\rule{#1}{0.4pt}}

% ===== Boîtes =====
\newtcolorbox{objectifbox}[1][]{colback=blue!5,colframe=blue!50,
    fonttitle=\bfseries,title=#1,breakable}
\newtcolorbox{infobox}[1][]{colback=green!5,colframe=green!50!black,
    fonttitle=\bfseries,title=#1,breakable}
\newtcolorbox{attentionbox}[1][]{colback=orange!5,colframe=orange!70,
    fonttitle=\bfseries,title=#1,breakable}
\newtcolorbox{retenirbox}[1][]{colback=yellow!10,colframe=yellow!50!black,
    fonttitle=\bfseries,title=#1,breakable}
\newtcolorbox{prereqbox}[1][]{colback=cyan!5,colframe=cyan!50!black,
    fonttitle=\bfseries,title=#1,breakable}
\newtcolorbox{theorembox}[1][]{colback=purple!5,colframe=purple!50,
    fonttitle=\bfseries,title=#1,breakable}
\newtcolorbox{appliboxQ}[1][]{colback=gray!8,colframe=gray!60,
    fonttitle=\bfseries,title=#1,breakable}
\newtcolorbox{appliboxR}[1][]{colback=green!3,colframe=green!40!black,
    fonttitle=\bfseries\itshape,title=#1,breakable}
\newtcolorbox{exercicebox}[1][]{colback=white,colframe=gray!70,
    fonttitle=\bfseries,title=#1,breakable}
\newtcolorbox{lienbox}[1][]{colback=blue!3,colframe=blue!70!black,
    fonttitle=\bfseries\itshape,title=#1,breakable,
    before upper={\itshape}}

% ===== Code =====
\lstset{
  language=Java,
  basicstyle=\ttfamily\small,
  backgroundcolor=\color{gray!8},
  frame=single,
  framerule=0.5pt,
  rulecolor=\color{gray!40},
  keywordstyle=\color{blue!70!black}\bfseries,
  commentstyle=\color{green!50!black}\itshape,
  stringstyle=\color{red!60!black},
  numbers=left,
  numberstyle=\tiny\color{gray},
  breaklines=true,
  breakatwhitespace=true,
  showstringspaces=false,
  tabsize=4,
  extendedchars=true,
  xleftmargin=4pt,
  xrightmargin=4pt,
  aboveskip=8pt,
  belowskip=8pt,
  literate={é}{{\'e}}1 {è}{{\`e}}1 {ê}{{\^e}}1 {à}{{\`a}}1 {ù}{{\`u}}1
           {ô}{{\^o}}1 {î}{{\^i}}1 {ç}{{\c{c}}}1 {û}{{\^u}}1 {â}{{\^a}}1
           {ë}{{\"e}}1 {ï}{{\"i}}1 {É}{{\'E}}1 {È}{{\`E}}1 {À}{{\`A}}1
           {Ç}{{\c{C}}}1 {Ê}{{\^E}}1 {œ}{{\oe}}1 {°}{{$^\circ$}}1
           {«}{{\guillemotleft}}1 {»}{{\guillemotright}}1
}
\lstdefinestyle{bash}{language=bash,keywordstyle=\color{black},
  commentstyle=\color{gray!70!black}\itshape,numbers=none}

% ===== En-têtes =====
\pagestyle{fancy}
\fancyhf{}
\fancyhead[L]{\small Lab S01 : ADTs et Interfaces Java}
\fancyhead[R]{\small ESP/UCAD --- M1}
\fancyfoot[L]{\small Dr. El Hadji Bassirou TOURÉ}
\fancyfoot[C]{\thepage}
\fancyfoot[R]{\small 2025--2026}
\renewcommand{\headrulewidth}{0.4pt}
\renewcommand{\footrulewidth}{0.4pt}

\title{%
\vspace{-1cm}
\textbf{\LARGE Structures de Données et}\\
\textbf{\LARGE Algorithmes Avancés}\\[0.8cm]
\textbf{\huge Lab S01 --- ADTs et Interfaces Java}\\[0.3cm]
\textbf{\Large Implémenter \texttt{ArrayDeque<T>} et le mettre à l'épreuve}\\[0.5cm]
\large Master 1 --- 2025--2026}
\author{Dr. El Hadji Bassirou TOURÉ\\
Département Génie Informatique\\
Université Cheikh Anta Diop de Dakar}
\date{}

\begin{document}
\maketitle
\thispagestyle{empty}

\begin{prereqbox}[Avant de commencer]
\textbf{Prérequis :} TP-1 complété (Docker, Gradle, JUnit opérationnels) et
CM~S01 lu.\\[4pt]
\textbf{Organisation :} ce document comporte deux parties.
\emph{Partie I} --- vous implémentez et vous observez.
\emph{Partie II} --- vous établissez par le calcul ce que vous venez
d'observer.\\[4pt]
\textbf{Ce Lab n'est pas noté.} Il est formatif : son rôle est de rendre
exécutable ce que le CM a énoncé.\\[4pt]
\textbf{Livrable :} le dépôt Git du projet \texttt{lab-s01} et le rapport
individuel décrit en fin de document.
\end{prereqbox}

\begin{objectifbox}[Ce que vous saurez faire à la fin de ce Lab]
\begin{enumerate}[nosep]
    \item Créer un projet Gradle et y déclarer une interface générique.
    \item Implémenter \texttt{ArrayDeque<T>} sur un tableau circulaire :
          les quatre opérations d'extrémité et le redimensionnement.
    \item Écrire un test JUnit qui \emph{force} le rebouclage du tableau,
          et un autre qui force le redimensionnement.
    \item Lire un programme d'observation et y reconnaître le déplacement
          des indices \texttt{tete} et \texttt{fin}.
    \item Restreindre un deque pour en faire une pile, et évaluer une
          expression postfixe.
    \item Calculer à la main l'état d'un tableau circulaire après une
          séquence d'opérations.
\end{enumerate}
\end{objectifbox}

\tableofcontents
\newpage

% #####################################################################
\part{Pratique : implémentation et observation}
% #####################################################################

\section{Environnement de travail}

Le travail se fait dans le conteneur installé au TP-1 : vous éditez les
fichiers avec votre éditeur habituel sur votre machine, vous compilez et
testez dans le conteneur.

\begin{lstlisting}[style=bash,caption={Rappel : lancer le conteneur depuis le dossier de travail}]
# Windows (Git Bash)
MSYS_NO_PATHCONV=1 docker run -it -v "$(pwd):/workspace" elbachirtoure/dsa-java-m1:v1.1

# macOS / Linux
docker run -it -v "$(pwd):/workspace" elbachirtoure/dsa-java-m1:v1.1
\end{lstlisting}

\begin{lstlisting}[style=bash,caption={Créer le projet}]
cd /workspace
new-project lab-s01
cd lab-s01
gradle build --no-daemon
\end{lstlisting}

\begin{attentionbox}[Deux règles héritées du TP-1]
La classe principale doit s'appeler \texttt{Main} --- c'est ce que
\texttt{new-project} inscrit dans \texttt{build.gradle}. Et toutes les
commandes Gradle prennent \texttt{--no-daemon} : sans cette option, le
démon Gradle échoue sur les volumes montés.
\end{attentionbox}

% =====================================================================
\section{Étape 1 --- Poser le contrat}

\begin{lienbox}[Ce qu'on matérialise du CM]
L'ADT est un contrat, et en Java un contrat s'écrit \texttt{interface}. Rien
dans ce fichier ne dit comment les éléments seront rangés : c'est
précisément ce qui permettra d'en écrire plusieurs implémentations.
\end{lienbox}

Créez \texttt{src/main/java/Deque.java} :

\begin{lstlisting}[caption={\texttt{Deque.java} --- le contrat, fourni}]
public interface Deque<T> {
    void addFirst(T element);
    void addLast(T element);
    T removeFirst();
    T removeLast();
    T peekFirst();
    T peekLast();
    int size();
    boolean isEmpty();
}
\end{lstlisting}

% =====================================================================
\section{Étape 2 --- Le squelette et les champs}

\begin{lienbox}[Ce qu'on matérialise du CM]
Le tableau circulaire ne déplace jamais les éléments : il déplace les
indices. Trois champs suffisent --- le tableau, l'indice \texttt{tete} du
premier élément, et le nombre d'éléments \texttt{taille}. L'indice de la
prochaine case libre en fin se calcule par $(\texttt{tete} + \texttt{taille})
\bmod m$.
\end{lienbox}

Créez \texttt{src/main/java/ArrayDeque.java} :

\begin{lstlisting}[caption={\texttt{ArrayDeque.java} --- champs et utilitaires}]
public class ArrayDeque<T> implements Deque<T> {

    private Object[] donnees;
    private int tete;      // indice du premier élément
    private int taille;    // nombre d'éléments présents

    public ArrayDeque() {
        this.donnees = new Object[8];
        this.tete = 0;
        this.taille = 0;
    }

    public int size()          { return taille; }
    public boolean isEmpty()   { return taille == 0; }
    public int capacite()      { return donnees.length; }  // utile pour observer
    public int indiceTete()    { return tete; }            // utile pour observer

    /** Avance ou recule un indice sur l'anneau. */
    private int suivant(int i) { return (i + 1) % donnees.length; }
    private int precedent(int i) { return (i - 1 + donnees.length) % donnees.length; }
}
\end{lstlisting}

\begin{attentionbox}[Le point sur lequel presque tout le monde trébuche]
Dans \texttt{precedent}, le terme \texttt{+ donnees.length} n'est pas
décoratif. En Java, \texttt{(0 - 1) \% 8} vaut $-1$, pas $7$ : l'opérateur
\texttt{\%} renvoie un reste dont le signe suit le dividende. Sans ce
correctif, le premier \texttt{addFirst} sur un deque vide lève une
\texttt{ArrayIndexOutOfBoundsException}.
\end{attentionbox}

% =====================================================================
\section{Étape 3 --- Les quatre opérations d'extrémité}

\begin{lienbox}[Ce qu'on matérialise du CM]
Les quatre opérations écrivent ou lisent une seule case et bougent un seul
indice. Aucune boucle : c'est la raison pour laquelle le deque est la
structure de choix quand on travaille aux deux bouts.
\end{lienbox}

\begin{lstlisting}[caption={Les quatre opérations}]
    public void addLast(T element) {
        if (taille == donnees.length) agrandir();
        donnees[(tete + taille) % donnees.length] = element;
        taille++;
    }

    public void addFirst(T element) {
        if (taille == donnees.length) agrandir();
        tete = precedent(tete);
        donnees[tete] = element;
        taille++;
    }

    @SuppressWarnings("unchecked")
    public T removeFirst() {
        if (isEmpty()) throw new java.util.NoSuchElementException("deque vide");
        T premier = (T) donnees[tete];
        donnees[tete] = null;
        tete = suivant(tete);
        taille--;
        return premier;
    }

    @SuppressWarnings("unchecked")
    public T removeLast() {
        if (isEmpty()) throw new java.util.NoSuchElementException("deque vide");
        int dernier = (tete + taille - 1) % donnees.length;
        T valeur = (T) donnees[dernier];
        donnees[dernier] = null;
        taille--;
        return valeur;
    }
\end{lstlisting}

\begin{exercicebox}[Vérification 1 --- à compléter]
Écrivez \texttt{peekFirst()} et \texttt{peekLast()} sur le modèle
ci-dessus. Ces deux méthodes ne modifient ni \texttt{tete} ni
\texttt{taille}.

Quelle expression donne l'indice du dernier élément ? \lig{7cm}
\end{exercicebox}

% =====================================================================
\section{Étape 4 --- Le redimensionnement}

\begin{lienbox}[Ce qu'on matérialise du CM]
La capacité d'un tableau est finie. Quand le deque est plein, on alloue un
tableau deux fois plus grand et on y recopie les éléments \emph{dans
l'ordre logique}, en repartant de \texttt{tete} : le nouveau tableau n'est
plus rebouclé, et \texttt{tete} redevient $0$.
\end{lienbox}

\begin{lstlisting}[caption={Agrandir en remettant l'anneau à plat}]
    private void agrandir() {
        Object[] neuf = new Object[donnees.length * 2];
        for (int k = 0; k < taille; k++) {
            neuf[k] = donnees[(tete + k) % donnees.length];  // ordre logique
        }
        donnees = neuf;
        tete = 0;
    }
\end{lstlisting}

\begin{attentionbox}[L'erreur à ne pas commettre]
Recopier avec \texttt{System.arraycopy(donnees, 0, neuf, 0, taille)} paraît
plus rapide, mais copie les cases dans l'ordre \emph{physique}. Si le deque
a rebouclé, l'ordre logique est détruit et les tests d'ordre échouent sans
que rien ne plante.
\end{attentionbox}

% =====================================================================
\section{Étape 5 --- Les tests}

\begin{lienbox}[Ce qu'on matérialise du CM]
Un test vérifie le \emph{contrat}, pas l'implémentation : il n'a pas à
connaître \texttt{tete}. Mais deux phénomènes internes doivent être
provoqués volontairement, car ils ne surviennent pas dans un usage naïf : le
rebouclage et le redimensionnement.
\end{lienbox}

Créez \texttt{src/test/java/ArrayDequeTest.java} :

\begin{lstlisting}[caption={Tests JUnit fournis}]
import org.junit.jupiter.api.Test;
import static org.junit.jupiter.api.Assertions.*;

public class ArrayDequeTest {

    @Test
    void ordreFifoEtLifo() {
        Deque<String> d = new ArrayDeque<>();
        d.addLast("Fatou");
        d.addLast("Moussa");
        assertEquals("Fatou", d.removeFirst());   // comportement de file
        d.addFirst("Awa");
        assertEquals("Awa", d.removeFirst());     // comportement de pile
        assertEquals(1, d.size());
    }

    @Test
    void testForcerLeRebouclage() {
        // capacité initiale 8 : on remplit, on vide par devant, on remplit encore
        Deque<Integer> d = new ArrayDeque<>();
        for (int i = 0; i < 6; i++) d.addLast(i);
        for (int i = 0; i < 4; i++) d.removeFirst();   // tete avance à 4
        for (int i = 100; i < 104; i++) d.addLast(i);  // la fin repasse par 0
        assertEquals(6, d.size());
        assertEquals(4, d.removeFirst());
        assertEquals(103, d.removeLast());
    }

    @Test
    void testForcerLeRedimensionnement() {
        Deque<Integer> d = new ArrayDeque<>();
        for (int i = 0; i < 100; i++) d.addFirst(i);   // dépasse largement 8
        assertEquals(100, d.size());
        assertEquals(99, d.removeFirst());             // ordre préservé
        assertEquals(0, d.removeLast());
    }

    @Test
    void videLeveUneException() {
        Deque<String> d = new ArrayDeque<>();
        assertThrows(java.util.NoSuchElementException.class, d::removeFirst);
    }
}
\end{lstlisting}

\begin{lstlisting}[style=bash,caption={Lancer les tests}]
gradle test --no-daemon
\end{lstlisting}

\begin{exercicebox}[Vérification 2 --- à compléter]
Ajoutez deux tests de votre cru :
\begin{enumerate}[nosep]
    \item un test qui alterne \texttt{addFirst} et \texttt{addLast} puis
          vérifie l'ordre complet de sortie ;
    \item un test qui vide entièrement le deque puis le remplit à nouveau.
\end{enumerate}
Lequel des deux échoue si l'on oublie le \texttt{+ donnees.length} de
\texttt{precedent} ? \lig{6cm}
\end{exercicebox}

% =====================================================================
\section{Étape 6 --- Observer}

\begin{lienbox}[Ce qu'on matérialise du CM]
Le CM affirme que les éléments ne bougent jamais et que ce sont les indices
qui se déplacent. Ce programme le rend visible : après chaque opération, il
affiche \texttt{tete}, l'indice de fin, la taille et la capacité, et signale
les redimensionnements.
\end{lienbox}

\begin{lstlisting}[caption={\texttt{Main.java} --- programme d'observation}]
public class Main {

    static <T> void trace(String operation, ArrayDeque<T> d) {
        int fin = (d.indiceTete() + d.size()) % d.capacite();
        System.out.printf("%-16s tete=%2d  fin=%2d  taille=%2d  capacite=%2d%n",
                          operation, d.indiceTete(), fin, d.size(), d.capacite());
    }

    public static void main(String[] args) {
        ArrayDeque<String> d = new ArrayDeque<>();
        System.out.println("--- on remplit par la fin ---");
        for (String s : new String[]{"A", "B", "C"}) { d.addLast(s); trace("addLast(" + s + ")", d); }

        System.out.println("--- on remplit par le debut : tete recule ---");
        for (String s : new String[]{"D", "E"}) { d.addFirst(s); trace("addFirst(" + s + ")", d); }

        System.out.println("--- on retire aux deux bouts ---");
        d.removeLast();  trace("removeLast()", d);
        d.addFirst("F"); trace("addFirst(F)", d);
        d.removeFirst(); trace("removeFirst()", d);

        System.out.println("--- on depasse la capacite ---");
        for (int i = 0; i < 12; i++) { d.addLast("x" + i); trace("addLast(x" + i + ")", d); }
    }
}
\end{lstlisting}

\begin{lstlisting}[style=bash,caption={Exécuter}]
gradle run --no-daemon
\end{lstlisting}

\begin{exercicebox}[Vérification 3 --- notez vos observations]
\begin{enumerate}[nosep]
    \item Après le premier \texttt{addFirst}, que vaut \texttt{tete} ?
          \lig{2.5cm} \quad Pourquoi cette valeur ? \lig{6cm}
    \item À quelle opération la capacité change-t-elle pour la première
          fois ? \lig{5cm}
    \item Juste après un changement de capacité, que vaut \texttt{tete} ?
          \lig{2.5cm} \quad Expliquez : \lig{7cm}
    \item Relevez une ligne où \texttt{fin} est \emph{inférieur} à
          \texttt{tete}. Que cela signifie-t-il ? \lig{7cm}
\end{enumerate}
\end{exercicebox}

% =====================================================================
\section{Étape 7 --- Une pile, et un évaluateur postfixe}

\begin{lienbox}[Ce qu'on matérialise du CM]
Le deque subsume la pile : en n'utilisant que \texttt{addLast} et
\texttt{removeLast}, on obtient exactement le comportement
dernier-arrivé-premier-sorti. On n'écrit donc pas une nouvelle structure ---
on restreint celle qu'on a.
\end{lienbox}

\begin{lstlisting}[caption={\texttt{PostfixEvaluator.java}}]
public class PostfixEvaluator {

    public static int evaluer(String expression) {
        ArrayDeque<Integer> pile = new ArrayDeque<>();
        for (String jeton : expression.trim().split("\\s+")) {
            if (estOperateur(jeton)) {
                if (pile.size() < 2) throw new IllegalArgumentException("operandes manquants");
                int b = pile.removeLast();     // l'ordre compte
                int a = pile.removeLast();
                pile.addLast(appliquer(jeton, a, b));
            } else {
                pile.addLast(Integer.parseInt(jeton));
            }
        }
        int resultat = pile.removeLast();
        if (!pile.isEmpty()) throw new IllegalArgumentException("expression mal formee");
        return resultat;
    }

    private static boolean estOperateur(String j) {
        return j.equals("+") || j.equals("-") || j.equals("*") || j.equals("/");
    }

    private static int appliquer(String op, int a, int b) {
        switch (op) {
            case "+": return a + b;
            case "-": return a - b;
            case "*": return a * b;
            case "/": return a / b;
            default: throw new IllegalArgumentException("operateur inconnu : " + op);
        }
    }
}
\end{lstlisting}

\begin{exercicebox}[Vérification 4 --- à compléter]
Écrivez trois tests JUnit pour \texttt{PostfixEvaluator} :
\begin{enumerate}[nosep]
    \item \texttt{"6 2 / 3 4 * +"} doit valoir \lig{1.5cm}
    \item une expression contenant une soustraction, pour détecter une
          inversion des opérandes ;
    \item une expression mal formée, qui doit lever une exception.
\end{enumerate}
\end{exercicebox}

% =====================================================================
\section{Enregistrer le travail}

\begin{lstlisting}[style=bash,caption={Git}]
cd /workspace/lab-s01
git init
git add .
git commit -m "Lab S01 : ArrayDeque, tests, evaluateur postfixe"
\end{lstlisting}

\newpage
% #####################################################################
\part{Théorie : exercices d'application}
% #####################################################################

\begin{infobox}[Comment aborder cette partie]
Ces exercices ne sont pas un contrôle et ne sont pas notés. Ils reprennent
exactement les situations que vous venez de programmer : le but est que vous
retrouviez par le calcul ce que le programme d'observation vous a montré à
l'écran. Remplissez directement sur ce document.
\end{infobox}

% ---------------------------------------------------------------------
\section{Exercice 1 --- Arithmétique de l'anneau}

Un \texttt{ArrayDeque} de capacité $m = 8$, initialement vide, avec
\texttt{tete} $= 0$. Complétez le tableau après chaque opération. Les cases
du tableau interne sont numérotées de $0$ à $7$.

\vspace{0.3cm}
\begin{center}
\begin{tabular}{lccc}
\toprule
\textbf{Opération} & \textbf{tete} & \textbf{taille} & \textbf{case écrite ou libérée} \\
\midrule
\texttt{addLast(A)}  & \rule{0pt}{0.75cm} & & \\
\texttt{addLast(B)}  & \rule{0pt}{0.75cm} & & \\
\texttt{addLast(C)}  & \rule{0pt}{0.75cm} & & \\
\texttt{addFirst(D)} & \rule{0pt}{0.75cm} & & \\
\texttt{addFirst(E)} & \rule{0pt}{0.75cm} & & \\
\texttt{removeLast()} & \rule{0pt}{0.75cm} & & \\
\texttt{addFirst(F)} & \rule{0pt}{0.75cm} & & \\
\texttt{removeFirst()} & \rule{0pt}{0.75cm} & & \\
\bottomrule
\end{tabular}
\end{center}

\vspace{0.3cm}
\textit{Aide : la case écrite par \texttt{addLast} est
$(\texttt{tete} + \texttt{taille}) \bmod 8$ ; \texttt{addFirst} écrit dans
$(\texttt{tete} - 1 + 8) \bmod 8$.}

\vspace{0.4cm}
\noindent\textbf{(a)} Quelles valeurs ont été retirées, et dans cet ordre ?
\lig{8cm}

\vspace{0.3cm}
\noindent\textbf{(b)} Dessinez l'état final des huit cases :

\begin{center}
\begin{tabular}{cccccccc}
\toprule
0 & 1 & 2 & 3 & 4 & 5 & 6 & 7 \\
\midrule
\rule{0pt}{0.9cm}\lig{1.2cm} & \lig{1.2cm} & \lig{1.2cm} & \lig{1.2cm} &
\lig{1.2cm} & \lig{1.2cm} & \lig{1.2cm} & \lig{1.2cm} \\
\bottomrule
\end{tabular}
\end{center}

\vspace{0.3cm}
\noindent\textbf{(c)} À la fin, \texttt{tete} vaut-il un indice inférieur ou
supérieur à celui de la case de fin ? Qu'est-ce que cela indique sur l'état
du tableau ? \lig{10cm}

% ---------------------------------------------------------------------
\section{Exercice 2 --- Le modulo négatif}

\noindent\textbf{(a)} En Java, donnez la valeur de chaque expression :

\vspace{0.2cm}
\begin{center}
\begin{tabular}{lc}
\toprule
\textbf{Expression} & \textbf{Valeur} \\
\midrule
\texttt{(3 - 1) \% 8}     & \rule{0pt}{0.7cm} \\
\texttt{(0 - 1) \% 8}     & \rule{0pt}{0.7cm} \\
\texttt{(0 - 1 + 8) \% 8} & \rule{0pt}{0.7cm} \\
\texttt{(7 + 1) \% 8}     & \rule{0pt}{0.7cm} \\
\bottomrule
\end{tabular}
\end{center}

\vspace{0.3cm}
\noindent\textbf{(b)} Un étudiant écrit \texttt{tete = (tete - 1) \% donnees.length;}
puis appelle \texttt{addFirst} sur un deque vide. Que se passe-t-il
exactement, et à quelle ligne ? \lig{10cm}

\vspace{0.2cm}
\noindent\lig{14cm}

\vspace{0.3cm}
\noindent\textbf{(c)} Ce bug est-il détecté par un test qui n'utilise que
\texttt{addLast} et \texttt{removeFirst} ? Justifiez. \lig{9cm}

% ---------------------------------------------------------------------
\section{Exercice 3 --- Trace d'une pile}

Déroulez la vérification de parenthésage sur
\texttt{\{ a ( b [ c ] ) d \}} en notant le contenu de la pile après chaque
symbole significatif (on ignore les lettres).

\vspace{0.3cm}
\begin{center}
\begin{tabular}{lll}
\toprule
\textbf{Symbole lu} & \textbf{Action} & \textbf{Pile après} \\
\midrule
\rule{0pt}{0.7cm} & & \\
\rule{0pt}{0.7cm} & & \\
\rule{0pt}{0.7cm} & & \\
\rule{0pt}{0.7cm} & & \\
\rule{0pt}{0.7cm} & & \\
\rule{0pt}{0.7cm} & & \\
\bottomrule
\end{tabular}
\end{center}

\vspace{0.3cm}
\noindent\textbf{(a)} L'expression est-elle équilibrée ? \lig{3cm}

\vspace{0.2cm}
\noindent\textbf{(b)} Modifiez un seul symbole de l'expression pour qu'elle
soit rejetée \emph{par le test final} (pile non vide à la fin), et non par
une correspondance fautive : \lig{7cm}

% ---------------------------------------------------------------------
\section{Exercice 4 --- Trace postfixe}

Évaluez \texttt{6 2 / 3 4 * + 5 -} en notant la pile après chaque jeton.

\vspace{0.3cm}
\begin{center}
\begin{tabular}{ccccccccc}
\toprule
\texttt{6} & \texttt{2} & \texttt{/} & \texttt{3} & \texttt{4} &
\texttt{*} & \texttt{+} & \texttt{5} & \texttt{-} \\
\midrule
\rule{0pt}{1.8cm} & & & & & & & & \\
\bottomrule
\end{tabular}
\end{center}

\vspace{0.3cm}
\noindent\textbf{(a)} Résultat : \lig{2cm}

\vspace{0.2cm}
\noindent\textbf{(b)} Hauteur maximale atteinte par la pile : \lig{2cm}

\vspace{0.2cm}
\noindent\textbf{(c)} Si l'on inverse les deux dépilements dans le code, quel
résultat obtient-on pour cette expression ? \lig{4cm}

% ---------------------------------------------------------------------
\section{Exercice 5 --- Contrat et coût}

Pour chaque opération, indiquez si le travail effectué est indépendant du
nombre $n$ d'éléments présents, ou s'il croît avec $n$.

\vspace{0.3cm}
\begin{center}
\begin{tabular}{lcc}
\toprule
\textbf{Opération} & \textbf{Tableau contigu} & \textbf{Nœuds chaînés} \\
\midrule
\texttt{get(i)}                & \rule{0pt}{0.7cm} & \\
\texttt{ajouter(x)} en fin     & \rule{0pt}{0.7cm} & \\
\texttt{ajouter(0, x)} en tête & \rule{0pt}{0.7cm} & \\
\texttt{supprimer(0)}          & \rule{0pt}{0.7cm} & \\
\texttt{taille()}              & \rule{0pt}{0.7cm} & \\
\bottomrule
\end{tabular}
\end{center}

\vspace{0.3cm}
\noindent\textbf{(a)} Un collègue implémente une file avec une
\texttt{ArrayList} : \texttt{enqueue} $\rightarrow$ \texttt{add(x)},
\texttt{dequeue} $\rightarrow$ \texttt{remove(0)}. Le contrat de l'ADT Queue
est-il respecté ? \lig{4cm}

\vspace{0.2cm}
\noindent\textbf{(b)} Un test JUnit portant uniquement sur le contrat
peut-il révéler le problème ? \lig{8cm}

\vspace{0.2cm}
\noindent\textbf{(c)} Comment le mettriez-vous en évidence expérimentalement ?
\lig{11cm}

\vspace{0.2cm}
\noindent\lig{14cm}

% ---------------------------------------------------------------------
\section{Exercice 6 --- Question subsidiaire}

On voudrait ajouter au contrat \texttt{Deque<T>} une opération
\texttt{get(int i)} renvoyant le $i$-ième élément depuis la tête.

\vspace{0.2cm}
\noindent\textbf{(a)} Écrivez l'expression de l'indice interne correspondant :
\lig{7cm}

\vspace{0.2cm}
\noindent\textbf{(b)} Cette opération a-t-elle sa place dans le contrat d'un
deque ? Argumentez en une ou deux phrases. \lig{10cm}

\vspace{0.2cm}
\noindent\lig{14cm}

\newpage
% =====================================================================
\section*{Pour aller plus loin (optionnel)}
\addcontentsline{toc}{section}{Pour aller plus loin (optionnel)}

\begin{itemize}[nosep]
    \item Écrire \texttt{LinkedDeque<T>}, seconde implémentation du même
          contrat, à l'aide de nœuds doublement chaînés --- et faire passer
          \emph{les mêmes tests} sans y changer une ligne. C'est la
          démonstration expérimentale du principe de la séance.
    \item Ajouter au deque une capacité maximale avec éviction automatique
          du plus ancien : vous obtenez un historique borné.
    \item Étendre l'évaluateur postfixe à la puissance et aux nombres
          négatifs.
\end{itemize}

% =====================================================================
\section*{Aide-mémoire}
\addcontentsline{toc}{section}{Aide-mémoire}

\begin{center}
\begin{tabular}{ll}
\toprule
\textbf{Question} & \textbf{Réponse} \\
\midrule
Indice de la prochaine case en fin & $(\texttt{tete} + \texttt{taille}) \bmod m$ \\
Indice du dernier élément & $(\texttt{tete} + \texttt{taille} - 1) \bmod m$ \\
Reculer un indice & $(i - 1 + m) \bmod m$ \\
Avancer un indice & $(i + 1) \bmod m$ \\
Deque vide & \texttt{taille == 0} (jamais \texttt{tete == fin}) \\
Après \texttt{agrandir()} & \texttt{tete} redevient $0$ \\
\midrule
Lancer le conteneur & \texttt{docker run -it -v "\$(pwd):/workspace" ...:v1.1} \\
Créer un projet & \texttt{new-project lab-s01} \\
Compiler & \texttt{gradle build -\/-no-daemon} \\
Tester & \texttt{gradle test -\/-no-daemon} \\
Exécuter & \texttt{gradle run -\/-no-daemon} \\
\bottomrule
\end{tabular}
\end{center}

% =====================================================================
\begin{attentionbox}[Rapport individuel --- obligatoire]
Chaque étudiant rédige un rapport de \textbf{1 à 2 pages} :
\begin{enumerate}[nosep]
    \item \textbf{Ce que j'ai fait :} classes implémentées, tests écrits, ce
          qu'affiche \texttt{Main}.
    \item \textbf{Les observations clés :} qu'est-ce que le programme
          d'observation m'a montré que je n'avais pas bien compris du CM ?
    \item \textbf{Les difficultés :} erreurs de compilation, bugs, blocages
          sur les exercices théoriques.
    \item \textbf{Ce que j'ai appris :} le concept le plus important retenu
          de cette séance.
\end{enumerate}
\textbf{Format :} PDF. \textbf{Nom :} \texttt{rapport-lab-S01-PRENOM-NOM.pdf}
\end{attentionbox}

\end{document}
